% ============================================================================
% paper.tex — QF (rQUF2e) build of the cluster-lineage manuscript.
% AI-DRAFTED (Claude), 2026-08-17, produced from revision 2 of
% cluster_lineage_manuscript.tex at owner instruction. THIS IS NOW THE LIVE
% MANUSCRIPT FILE; cluster_lineage_manuscript.tex is the frozen revision-2
% archive (one-canonical-version hygiene rule).
% Class mechanics per rQUFguide.tex (v2.1) and the QF Instructions for
% Authors (updated 17 July 2026), .private/qf_author_instructions.txt:
%   - rQUF2e class, no size/geometry options, single column; rQUF.bst
%     author-date references (\citep/\citet unchanged).
%   - Frontmatter: \author with \affil and \received; abstract (unstructured,
%     200 words); keywords environment (4-6 terms); classcode (JEL).
%   - Tables: caption ABOVE via \tbl{...}{...}, rules \toprule/\colrule/
%     \botrule, \label before \end{table}, notes via \tabnote.
%   - Appendices AFTER the reference list, opened by \appendices.
%   - Back matter before references: Acknowledgements, Disclosure statement,
%     Declaration of generative AI use (REQUIRED by the 2026 instructions),
%     Funding, Data availability statement.
%   - 35-page limit inclusive of references and footnotes; word count to be
%     stated at submission; submission fee USD 162; single-blind review (for
%     the anonymous PDF, move \maketitle between \title and \author per the
%     guide); LaTeX submissions upload the PDF plus zipped source files.
%   - Theorem-like environments keep the draft's global counters so Lemma 1 /
%     Propositions 1-2 / P1-P7 keep matching the E-series reports; the
%     typesetter renumbers to journal style at proof stage if required.
%   - rQUF2e.cls and rQUF.bst live in this folder for compilation and stay
%     untracked in git (third-party files, repo hygiene rule). The bundled
%     2014-vintage .sty files from rquflatex.zip are NOT copied: system
%     TeX Live versions are newer.
% Revision history carried over from cluster_lineage_manuscript.tex:
% Full-manuscript revision of cluster_distribution_theory.tex (2026-08-13),
% AI-DRAFTED (Claude), revision 2, 2026-08-17, at owner instruction and per
% the owner rulings of 2026-08-17 (Gate 0 + phases + two-part empirical split).
% Revision 2 (2026-08-17, QF-baseline restructure per
% agents/2026-08-17_claude_qf_structure_review.md):
%   E2: 11 sections consolidated to 8 (applications merged; robustness and
%       limitations merged); all labels preserved.
%   E3/E10: Section 5 subsectioned (panels and data / scorecard / simulation).
%   E4: proofs moved to Appendix A; statements stay in the body.
%   E5: title shortened to the named object (Gate 0 confirmation pending).
%   E6: keywords and JEL added (owner to confirm).
%   E7: Taylor & Francis back-matter skeleton added.
%   E8: two new introduction paragraphs (QF lineage; stability literature),
%       both %% TOPIC-CANDIDATE; four verified references added to the bib
%       (onnela2003dynamics, tumminello2005tool, marti2021review,
%       raffinot2017hierarchical).
%   E12: AUM spelled out in prose; MCF struck from the acronym audit note.
%   Verbatim-reuse fix: the motivating-arithmetic sentence now appears once
%       (introduction); Section 3.2 closes with a back-reference instead.
%   Exhibit plan: figure/table TODOs now carry F1-F6 / TA-TF identifiers
%       mapping to agents/2026-08-17_claude_sol_todo_qf_baseline_exhibits.md.
% Revision 3 (2026-08-21, W3 integration per ROADMAP_manuscript_finalisation
% v3; all inserted numbers are gated F/G-stage values):
%   - U1 signal and risk rows re-based to the gated G0 headline-window
%     re-score (2026-08-21_sol_G0_report.md); the 2026-08-17 U1 signal
%     quotations are RETIRED (source artifact lost 2026-08-14, G0
%     reconciliation shows they are not reproducible).
%   - Signal claim re-framed: volatility reduction is the cross-universe
%     robust result (all four volatility CIs exclude zero; no return or
%     Sharpe CI does). F6 + G0 intervals quoted in the text.
%   - Calibration bridge table (TB) inserted from F1, with the
%     calibration-transfer disclosure (adopted 0.0866 vs own-cell 0.1753).
%   - Theory scorecard table (TD) inserted from F5; P2/P6/P7 verdict prose
%     added; P4-revised verdict from F2; simulation results from F4.
%   - Membership/interpretability numbers from F3; MCF paragraph drafted
%     from the F3 report (owner review under W1).
%   - Data availability statement replaced by the F10 draft (W2 sign-off).
%   - REMAINING [TODO]s: exhibit inserts from finalisation/exhibits (TA, TC,
%     TE/TF, figures F1-F6), U1 turnover cells and U1 budget-row re-scores
%     (trace via exhibit_index.csv / table_existing_*.csv), and the owner
%     items W1/W2 (title, mechanism passages, P4 proposition text,
%     sign-offs, notation table).
% W3 blocks are delimited by BEGIN/END marker comments where they replace
% or extend revision-2 text inside otherwise unchanged passages.
% Seam protocol: this whole file is one generated block, EXCEPT the central
% claim in the introduction, which is owner-drafted text (marked %% OWNER).
% Candidate mechanism sentences for owner rewrite are marked %% MECHANISM.
% Load-bearing topic sentences drafted ahead of the owner skeleton are
% marked %% TOPIC-CANDIDATE.
% Every pending analysis is marked [TODO: ...] and maps to the Part A / Part B
% items of agents/2026-08-17_claude_empirical_data_plan.md and the exhibit
% caps of agents/2026-08-17_claude_sol_todo_qf_baseline_exhibits.md.
% Number provenance: LaTeX comments name the source report in agents/ for
% every quoted number; the consolidated exhibit_index.csv (Phase C) will
% replace these comments as the traceability record.
% The theory sections carry over from cluster_distribution_theory.tex with
% integration edits only; that file remains the frozen 2026-08-13 archive.
% Compile: pdflatex + bibtex with cluster_lineage_refs.bib (rQUF.bst).
% The latexrelease rollback to 2018-12-01 is REQUIRED (rQUFguide.tex line 4):
% the 2014 class is incompatible with the current LaTeX kernel without it.
% booktabs must NOT be loaded: the class provides \toprule/\colrule/\botrule
% and clashes with booktabs' versions.
% ============================================================================
\RequirePackage[2018-12-01]{latexrelease}
\documentclass{rQUF2e}
\usepackage{amsthm}
\usepackage{url}

\theoremstyle{plain}
\newtheorem{lemma}{Lemma}
\newtheorem{proposition}{Proposition}
\newtheorem{corollary}{Corollary}
\newtheorem{assumption}{Assumption}
\newtheorem{prediction}{Prediction}
\theoremstyle{remark}
\newtheorem{remark}{Remark}

\newcommand{\E}{\mathbb{E}}
\newcommand{\Prob}{\mathbb{P}}
\newcommand{\R}{\mathbb{R}}
\newcommand{\vech}{\operatorname{vech}}
\newcommand{\ARI}{\operatorname{ARI}}
\newcommand{\VI}{\operatorname{VI}}
\newcommand{\TODO}[1]{\textbf{[TODO: #1]}}

% Working title, shortened to the named object per the QF-baseline review
% (titles in the target-journal sample run 4-17 words, median 8).
% [TODO: owner to confirm the title with the Gate 0 naming ruling.]
% For the anonymous review PDF, move \maketitle to directly after \title
% (rQUFguide section 3.1). Word count: [TODO: state at submission; the
% instructions require a word count with the paper.]
\title{Rolling-Ward clustering: noise-calibrated stability for rolling
correlation-based clusters}

\author{Artur SEPP$^{\ast}$\thanks{$^\ast$Corresponding author.
Email: \TODO{email}. \TODO{author list, full affiliations with postal
addresses (where the research was conducted), ORCiDs, and CRediT roles at
submission.} Reference implementation in the \texttt{factorlasso} package.}\\
\affil{\TODO{affiliation, full postal address}}
\received{working draft of 2026-08-17, revision 2; not yet submitted}}

\begin{document}
\maketitle

\begin{abstract}
Rolling correlation-based clustering supplies the peer groups behind risk
models, cross-sectional signals, and hierarchical allocation. Between two
monthly estimation dates a weekly EWMA estimator with span 156 replaces about
five percent of its weight, yet the unsmoothed Ward partition reassigns roughly
one quarter of assets, and the smoothing parameters that control this churn are
tuned rather than derived. We derive a distribution theory for the induced
partition process (stationarity, an exact increment law, and a boundary-flip
approximation) and we calibrate the smoothing threshold from the estimator's
own noise floor. On three estimation panels the calibrated threshold cuts
membership churn by up to 89\% at a bounded fidelity cost, and the theory
scorecard supports four predictions outright, two in revised form, and rejects
one. In three investment universes the stabilised clusters reduce the
volatility of matched momentum portfolios by 0.7 to 3.9 percentage points
against ranking controls, and every volatility interval excludes zero under a
block bootstrap while no return or Sharpe interval does. The smoothing
parameter thereby stops being a
tuning choice: a risk platform can set it from the span, the estimation
frequency, and a tail estimate, and defend the setting as a
statistical-evidence threshold.
\end{abstract}

\begin{keywords}
Correlation clustering; Hierarchical clustering; Cluster stability; EWMA;
Hierarchical risk parity; Cross-sectional momentum
% [TODO: owner to confirm the keyword set; the instructions require 4-6
% terms and this set carries 6.]
\end{keywords}

\begin{classcode}C38, C58, G11 \TODO{owner to confirm}\end{classcode}

% ----------------------------------------------------------------------------
\section{Introduction}
% ----------------------------------------------------------------------------

%% OWNER (central claim, drafted by the owner 2026-08-17; one correction
%% applied, "if" -> "is"; do not rewrite):
Estimator noise may elevate the membership churn in clusters produced from
rolling correlation. The noise is amplified at assignment boundaries. Deriving
the smoothing threshold from the estimator's noise floor stabilises clusters at
bounded fidelity cost. The stabilised clusters improve the risk properties of
cluster-scored signals and cluster-budgeted allocations.

%% TOPIC-CANDIDATE (QF-lineage paragraph; expanded in revision 2 per E8,
%% drafted ahead of the owner skeleton -- owner to rewrite topic and crux
%% sentences):
Correlation-based clustering of financial instruments is an established object
of study in quantitative finance. \citet{mantegna1999hierarchical} extracts a
hierarchical taxonomy of equities from the correlation matrix, and
\citet{onnela2003dynamics} track that taxonomy through rolling estimation
windows, documenting that the fitted trees evolve as the window moves.
\citet{tumminello2005tool} filter the correlation structure into sparse
graphs, and \citet{marti2021review} survey two decades of this
correlation-network literature. On the portfolio side, hierarchical risk
parity converts the fitted dendrogram into capital allocations
\citep{lopezdeprado2016building, lopezdeprado2020machine}, and
\citet{raffinot2017hierarchical} documents the out-of-sample appeal of
hierarchical allocation across clustering variants. This literature studies
the partition, or the portfolio built from it, at a point in time. The object
that a production risk platform consumes is the \emph{sequence} of partitions
re-estimated on a schedule from a rolling estimator, and the stability of this
sequence has no received calibration.

%% TOPIC-CANDIDATE (stability-literature paragraph; new in revision 2 per E8,
%% drafted ahead of the owner skeleton -- owner to rewrite topic and crux
%% sentences):
A second literature treats the stability of clustering directly, and it
supplies the constraint rather than the calibration. Stability alone is a
degenerate objective, because a frozen partition is perfectly stable
\citep{bendavid2006sober, vonluxburg2010clustering}. Evolutionary clustering
therefore trades a snapshot cost against a temporal cost, with a weight that
its users tune \citep{chakrabarti2006evolutionary, xu2014adaptive}, and
\citet{marti2016clustering} quantify how much sample the snapshot side of that
trade-off requires for financial correlations. We contribute the missing
piece: the temporal weight is derived, not tuned, from the sampling noise of
the EWMA correlation estimator itself.

The scale of the problem is visible in one piece of arithmetic. At weekly
frequency with span $N=156$ and monthly re-estimation, the EWMA estimator
replaces 5.4\% of its weight between estimation dates
(Section~\ref{sec:increments}). The unsmoothed Ward partition of our equity
panel nevertheless reassigns roughly 24\% of assets per month. A five percent
information innovation that produces a twenty-five percent membership
reassignment identifies the dendrogram cut as a noise amplifier, and the theory
below locates the amplification at the cluster boundaries.

We make four contributions. We derive the stationary law and the exact
increment law of the partition process induced by EWMA correlation estimation
(Section~\ref{sec:theory}). We prove a boundary-flip approximation for
membership churn under a block-structured population model and we obtain the
smoothing threshold as a designed hysteresis quantity calibrated from the
estimator's noise floor (Section~\ref{sec:smoothing}). We confirm the theory's
testable predictions on three estimation panels spanning equities, futures, and
multi-asset funds (Section~\ref{sec:confirmation}). We document, in three
investment universes, that the stabilised clusters improve the risk properties
of cluster-scored momentum signals and of cluster-budgeted long-only
allocations (Section~\ref{sec:applications}).

%% MECHANISM (candidate for owner rewrite): the substitution statement.
The calibration replaces two incumbents. Evolutionary clustering carries a
temporal-cost weight that its users tune \citep{chakrabarti2006evolutionary,
xu2014adaptive}. Production platforms instead freeze partitions between
scheduled reclustering dates chosen by convention. Our threshold substitutes
for both: it is derived from the span, the estimation frequency, and a tail
estimate, and switching costs one closed-form evaluation of
equation~\eqref{eq:deltalevel}. We name the underlying construction
\emph{Rolling-Ward} clustering, and we ship it, together with the smoothing
and lineage analytics, in the open-source \texttt{factorlasso}
package.\footnote{\url{https://github.com/ArturSepp/FactorLasso}.
\TODO{pin the release version that carries the paper's API at submission.}}

The paper proceeds as follows. Section~\ref{sec:rollingward} defines
Rolling-Ward clustering and its lineage overlay. Section~\ref{sec:theory}
derives the distribution theory, with proofs in
Appendix~\ref{app:proofs}. Section~\ref{sec:smoothing} states the smoothing
design problem, the noise-floor calibration, and the testable predictions.
Section~\ref{sec:confirmation} confirms the predictions on the three
estimation panels. Section~\ref{sec:applications} reports the two investment
applications: peer-contained momentum scoring and cluster-budgeted long-only
allocation. Section~\ref{sec:robustness} collects sensitivity material and
states limitations, and Section~\ref{sec:conclusion} concludes.

% ----------------------------------------------------------------------------
\section{Rolling-Ward clustering and lineage}
\label{sec:rollingward}
% ----------------------------------------------------------------------------

%% TOPIC-CANDIDATE
Rolling-Ward clustering is the composition of three standard components on a
rolling schedule. We observe returns $r_t \in \R^d$ on $d$ assets at a fixed
frequency. The covariance estimator is the exponentially weighted moving
average (EWMA) with span $N$,
\begin{equation}
\widehat{\Sigma}_t \;=\; \lambda\, \widehat{\Sigma}_{t-1}
\;+\; (1-\lambda)\, r_t r_t^{\top},
\qquad \lambda \;=\; 1 - \frac{2}{N+1},
\label{eq:ewma}
\end{equation}
with the correlation matrix
$\widehat{C}_t = \widehat{D}_t^{-1/2}\, \widehat{\Sigma}_t\, \widehat{D}_t^{-1/2}$,
where $\widehat{D}_t$ holds the diagonal of $\widehat{\Sigma}_t$
\citep{riskmetrics1996}. The clustering map takes the correlation matrix into a
partition of the assets,
\begin{equation}
Q_t \;=\; f\!\left(\widehat{C}_t;\, P\right),
\label{eq:clustermap}
\end{equation}
where $P$ collects the clustering parameters. In our implementation $f$ applies
Ward linkage \citep{ward1963hierarchical} to the distance matrix
$1-\widehat{C}_t$ and cuts the dendrogram at a fixed fraction of its maximal
merge height. The partition is re-estimated at scheduled dates spaced $k$
observation periods apart, on the point-in-time eligible asset set at each
date. The theory below holds $f$ measurable and treats the Ward specification
as one instance.

Two statistics measure partition changes. The \emph{churn} is the fraction of
assets whose cluster assignment changes between consecutive scheduled dates,
matched through cluster identities, and annualised to changes per asset-year.
The \emph{variation of information} is the metric
$\VI(Q, Q') = H(Q \mid Q') + H(Q' \mid Q)$ on partitions
\citep{meila2007comparing}. Churn is the practice-side statistic and $\VI$ is
the theory-side metric. Both appear in the empirical sections so that theory
and exhibits speak one language. The volatility dynamics of returns carry a
parameter set $\Theta$, for which we use GARCH-type specifications
\citep{bollerslev1986generalized}.

%% MECHANISM (candidate for owner rewrite): membership vs identity.
Adaptive membership and stable identity are separate objects. Membership is
adaptive: the current partition drives scores and budgets. Identity is a
reporting overlay: a min-cost-flow matcher links clusters across dates into
named tracks for lineage statistics and labelling, and it never feeds back into
estimation or portfolio construction. All lineage statistics in this paper come
from this overlay.

%% W3-BEGIN (MCF description drafted from the F3 report subsection,
%% agents/2026-08-20_sol_F3_report.md; owner review under W1):
The matcher builds a bipartite graph whose nodes are the fingerprinted
clusters at each estimation date. Edges connect clusters at consecutive dates
and across gaps of up to six dates, with gap edges discounted geometrically. A
deterministic maximum-weight matching selects a vertex-disjoint cover of
paths, and each predecessor-free path carries one persistent track identity.
The bridge edges let a track survive a temporary split or merge, because the
joint solve can route identity around the gap instead of handing it to the
locally best consecutive match. Unmatched qualifying alternatives annotate
splits and merges, and the remaining endpoints mark births and deaths.
%% W3-END

% ----------------------------------------------------------------------------
\section{A distribution theory for rolling partitions}
\label{sec:theory}
% ----------------------------------------------------------------------------

% Carried from cluster_distribution_theory.tex (2026-08-13) with integration
% edits only. Owner-gated verdicts of stage E3/E3b apply to Section 5.
% Revision 2: proofs moved to Appendix A per E4 (QF baseline: no in-text
% proof environments in any of the seven sample papers).

\subsection{The induced partition process and its stationary law}

%% MECHANISM (candidate for owner rewrite): the point of this section in one
%% sentence: the partition inherits a stationary law from the estimator for
%% free, so the scientific question is not whether the cluster distribution
%% exists but where it concentrates.
The recursion in equation~\eqref{eq:ewma} is a contraction, so the estimated
covariance forgets its initial condition geometrically. Existence of a
stationary cluster distribution is then an inheritance property, not a deep
fact. We record it as a lemma. Appendix~\ref{app:proofs} collects the proofs
of this section.

\begin{lemma}[Stationarity of the partition process]
\label{lem:stationarity}
Let $\{r_t\}$ be strictly stationary and ergodic with
$\E \log^{+} \| r_1 r_1^{\top} \| < \infty$. Then the recursion
\eqref{eq:ewma} admits the almost-surely convergent stationary solution
\begin{equation}
\widehat{\Sigma}_t \;=\; (1-\lambda) \sum_{j=0}^{\infty} \lambda^{j}\,
r_{t-j} r_{t-j}^{\top},
\label{eq:stationarysol}
\end{equation}
and $\{\widehat{\Sigma}_t\}$ is strictly stationary and ergodic. For any
measurable clustering map $f$, the partition process
$Q_t = f(\widehat{C}_t; P)$ is strictly stationary and ergodic on the finite
lattice of partitions of $d$ assets, and its stationary law is the pushforward
of the stationary law of $\widehat{C}_t$ under $f$.
\end{lemma}

\begin{remark}[Integrated GARCH]
\label{rem:igarch}
The EWMA filter is the integrated GARCH(1,1) filter \citep{riskmetrics1996}.
Under an integrated GARCH data-generating process, returns remain strictly
stationary while their unconditional variance is infinite
\citep{nelson1990stationarity}. Lemma~\ref{lem:stationarity} still applies,
because it requires only the logarithmic moment. Moment-based functionals of
the partition process are the objects that need care in this regime, and the
variance-based results below assume finite fourth moments.
\end{remark}

\begin{remark}[Where the stationary law concentrates]
\label{rem:concentration}
The content of the stationary law is its concentration, not its existence.
Under a block-structured population correlation matrix, exact recovery of the
block partition from $n$ observations requires the within-versus-between
separation to exceed a threshold of order $\sqrt{\log d / n}$
\citep{bunea2020model}. The EWMA estimator carries an effective sample size of
order $N$ (Proposition~\ref{prop:increments}), which does not grow with $d$,
and the random-matrix regime of exponentially weighted estimators keeps the
ratio $d/N$ finite \citep{pafka2004exponential}. The stationary law of $Q_t$
therefore concentrates on the population partition when separation clears the
noise floor and stays diffuse below it. We use this dichotomy as a working
frame and defer its full characterisation, including the removal of the market
mode before clustering \citep{laloux1999noise, lopezdeprado2020machine}, to a
companion study.
\end{remark}

\subsection{The law of covariance and correlation increments}
\label{sec:increments}

%% MECHANISM (candidate for owner rewrite): between two monthly estimation
%% dates the EWMA replaces only a few percent of its weight with new data, so
%% the correlation innovations are small and asymptotically Gaussian, and the
%% relevant scale is the estimator's own sampling noise.
The one-step increment of the estimator is exact and already informative. From
equation~\eqref{eq:ewma},
\begin{equation}
\widehat{\Sigma}_t - \widehat{\Sigma}_{t-1}
\;=\; (1-\lambda)\left( r_t r_t^{\top} - \widehat{\Sigma}_{t-1} \right),
\label{eq:increment}
\end{equation}
so the estimator mean-reverts toward the conditional second moment at rate
$1-\lambda$, with an innovation driven by the outer product $r_t r_t^{\top}$.
The next proposition collects the fluctuation results that calibrate all later
formulas.

\begin{proposition}[Increment law and noise floor]
\label{prop:increments}
Let returns be i.i.d.\ with finite fourth moments, let
$V = \operatorname{Var}\!\left(\vech(r_1 r_1^{\top})\right)$, and let
$\lambda = 1 - 2/(N+1)$. Then the following hold.
\begin{enumerate}
\item[(i)] The stationary process $\vech(\widehat{\Sigma}_t)$ has
autocovariance function $\gamma(h) = \lambda^{|h|} \gamma(0)$ with
\begin{equation}
\gamma(0) \;=\; \frac{1-\lambda}{1+\lambda}\, V \;=\; \frac{V}{N},
\label{eq:statvar}
\end{equation}
so the EWMA carries the effective sample size $N$ exactly.
\item[(ii)] The increment over $k$ periods satisfies
\begin{equation}
\operatorname{Var}\!\left( \vech(\widehat{\Sigma}_t)
- \vech(\widehat{\Sigma}_{t-k}) \right)
\;=\; 2\left(1 - \lambda^{k}\right) \gamma(0).
\label{eq:kstep}
\end{equation}
\item[(iii)] For elliptically distributed returns with kurtosis parameter
$\kappa$, the delta method applied to the correlation coefficient gives the
stationary standard error
\begin{equation}
\operatorname{se}\!\left(\widehat{\rho}_{ij}\right)
\;\approx\; \sqrt{1+\kappa}\; \frac{1-\rho_{ij}^2}{\sqrt{N}},
\label{eq:noisefloor}
\end{equation}
with $\kappa = 0$ in the Gaussian case
\citep{neudecker1990asymptotic}.
\item[(iv)] For a GARCH(1,1) data-generating process with Gaussian innovations
and parameters $\Theta = (\alpha, \beta)$ satisfying
$(\alpha+\beta)^2 + 2\alpha^2 < 1$, the unconditional excess-kurtosis parameter
is
\begin{equation}
\kappa_{\Theta} \;=\; \frac{2\alpha^2}{1 - (\alpha+\beta)^2 - 2\alpha^2},
\label{eq:garchkappa}
\end{equation}
so the volatility dynamics enter the noise floor \eqref{eq:noisefloor} through
one closed-form multiplier \citep{bollerslev1986generalized}.
\end{enumerate}
\end{proposition}

The distance entries $d_{ij} = 1 - \widehat{\rho}_{ij}$ inherit the standard
error \eqref{eq:noisefloor}. Combining parts (ii) and (iii) gives the
per-transition distance noise at estimation step $k$,
\begin{equation}
\sigma_d(k) \;=\; \sqrt{2\left(1-\lambda^{k}\right)}\;
\sqrt{1+\kappa}\; \frac{1-\bar{\rho}^2}{\sqrt{N}},
\label{eq:transitionnoise}
\end{equation}
evaluated at a representative within-cluster correlation $\bar{\rho}$. Under
misspecification, the same objects follow from the asymptotic filtering theory
of EWMA-type variance estimators \citep{nelson1992filtering,
nelson1994asymptotic}.
% [TODO: check the multivariate filtering reference (J. Econometrics 71, 1996)
% before citing it here.]

Equation \eqref{eq:transitionnoise} quantifies the motivating arithmetic of
the introduction. At weekly frequency with span $N=156$ and monthly estimation
($k \approx 4.3$), the factor $1-\lambda^{k}$ equals $0.054$, against a
baseline empirical churn of roughly one quarter of assets per month. The
amplification from a five percent innovation to a twenty-five percent
reassignment is what the next section locates at the cluster boundaries.

\subsection{Partition changes: a boundary-flip approximation}

%% MECHANISM (candidate for owner rewrite): the clustering map is flat almost
%% everywhere and jumps at assignment boundaries, so all churn comes from
%% assets whose margin is small relative to the estimator noise; smoothing is
%% hysteresis that widens the margin.
The clustering map is piecewise constant in the correlation matrix, so no delta
method applies to $\Delta Q_t$: the map is non-differentiable exactly at the
configurations that generate churn. The tractable route fixes a separated
population structure and studies which assets sit close to an assignment
boundary.

\begin{assumption}[Separated block structure]
\label{ass:block}
The population correlation matrix $C^{*}$ has a block structure with $G$
groups, in the sense of the $G$-block covariance models of
\citet{bunea2020model}, and the population dendrogram cut recovers the blocks.
For asset $i$, the population margin
\begin{equation}
m_i \;=\; \bar{d}\left(i, \mathcal{C}_{\mathrm{alt}(i)}\right)
- \bar{d}\left(i, \mathcal{C}(i)\right)
\;>\; 0
\label{eq:margin}
\end{equation}
is the gap between the mean distance from $i$ to its nearest alternative block
and the mean distance to its own block.
\end{assumption}

The margin in \eqref{eq:margin} uses mean distances, which is the
average-linkage functional. Ward linkage does not admit a local margin of this
form, because its merge criterion is a global functional of the tree. We
therefore prove the flip approximation for the flat cut under
Assumption~\ref{ass:block} and verify it for Ward by simulation
(Section~\ref{sec:confirmation}). This division is honest rather than
convenient: among linkage methods only single linkage is stable in the
Gromov--Hausdorff sense \citep{carlsson2010characterization}, so a
distributional theory for Ward partitions must run through a generative model,
not through generic metric stability.

\begin{proposition}[Boundary-flip approximation]
\label{prop:flip}
Under Assumption~\ref{ass:block}, let the estimated margin statistic
$\widehat{m}_{i,t}$ be asymptotically Gaussian around $m_i$ with standard
deviation $\sigma_i$ proportional to the per-transition distance noise
$\sigma_d(k)$ of equation \eqref{eq:transitionnoise}, with a proportionality
constant absorbing cluster sizes and the correlation of estimation errors.
Under a hysteresis discount $\delta \geq 0$ that retains asset $i$ in its
incumbent cluster unless the estimated margin moves against it by more than
$\delta$, the per-transition flip probability satisfies
\begin{equation}
\Prob\!\left(\text{flip}_i\right)
\;\approx\; \Phi\!\left( - \frac{m_i + \delta}{\sqrt{2}\,\sigma_i} \right),
\label{eq:flipprob}
\end{equation}
and the expected churn per asset-year is
\begin{equation}
\E\!\left[\text{churn}\right]
\;\approx\; \frac{A}{d} \sum_{i=1}^{d}
\Phi\!\left( - \frac{m_i + \delta}{\sqrt{2}\,\sigma_i} \right),
\label{eq:expectedchurn}
\end{equation}
where $A$ is the number of estimation dates per year.
\end{proposition}

% [TODO: OWNER (W2): restate the hysteresis component with the per-panel
% proportionality constant absorbed into sigma_i (owner ruling 10). The F2
% re-evaluation is COMPLETE (agents/2026-08-20_sol_F2_report.md): with c
% calibrated once on the baseline configuration, the revised prediction is
% supported as a cross-configuration ordering (correlations 0.992 equity
% full, 0.993 equity headline, 0.939 futures; fund rows descriptive), not as
% an equality in levels. Only the proposition text awaits the owner.]

% ----------------------------------------------------------------------------
\section{The smoothing design problem and its calibration}
\label{sec:smoothing}
% ----------------------------------------------------------------------------

Minimising partition variability alone is degenerate: any constant partition
achieves zero churn, so an unconstrained stability objective selects a frozen
or trivial clustering. This degeneracy is the central caveat of the clustering
stability literature, where stability measures boundary mass rather than
correctness \citep{bendavid2006sober, vonluxburg2010clustering}. We therefore
state the design problem as constrained:
\begin{equation}
\min_{P} \;\; \E\!\left[\text{churn}(P)\right]
\qquad \text{subject to} \qquad
\mathcal{F}(P) \;\geq\; \mathcal{F}_0,
\label{eq:designproblem}
\end{equation}
where $\mathcal{F}$ is a fidelity functional. In the theory, $\mathcal{F}$ is
the probability of recovering the population partition under
Assumption~\ref{ass:block}. In the implementation, $\mathcal{F}$ is the
adjusted Rand index band against the unsmoothed baseline together with a
resolution guard on the cluster-count distribution.

\begin{remark}[Relation to evolutionary clustering]
\label{rem:evolutionary}
The Lagrangian form of \eqref{eq:designproblem} is the snapshot-cost plus
temporal-cost objective of evolutionary clustering
\citep{chakrabarti2006evolutionary}, and similarity smoothing corresponds to
the adaptive forgetting estimators of that line \citep{xu2014adaptive}. Our
contribution at this point is not the objective but its calibration: the
temporal weight is derived from the estimator's noise floor through
Corollary~\ref{cor:delta} rather than tuned. Sample-length requirements for
correlation-based clustering in finance are studied by
\citet{marti2016clustering}, whose rates concern the time dimension of the
same trade-off.
\end{remark}

%% MECHANISM (candidate for owner rewrite): the design principle in one line:
%% a membership change should require statistical evidence that exceeds the
%% estimator's own sampling noise, which is the same gate logic as the
%% sign-constraint gate of our companion estimator.
Equation \eqref{eq:expectedchurn} turns the smoothing parameter into a designed
quantity. Two calibrations follow, and they differ in an observable way.

\begin{corollary}[Noise-floor calibration of the partition bonus]
\label{cor:delta}
Fix a false-flip budget $\alpha$ for zero-margin assets, with Gaussian quantile
$z_{1-\alpha}$. The level calibration sets the bonus against the stationary
noise of the distance level,
\begin{equation}
\delta^{*}_{\mathrm{lvl}}
\;=\; z_{1-\alpha}\, \sqrt{1+\kappa}\; \frac{1-\bar{\rho}^{\,2}}{\sqrt{N}},
\label{eq:deltalevel}
\end{equation}
and the innovation calibration sets it against the per-transition noise,
\begin{equation}
\delta^{*}_{\mathrm{inn}}
\;=\; \sqrt{2\left(1-\lambda^{k}\right)}\;\, \delta^{*}_{\mathrm{lvl}}.
\label{eq:deltainnovation}
\end{equation}
The level form does not depend on the estimation step $k$, while the
innovation form shrinks as estimation dates move closer together. The
dependence on $k$ therefore discriminates between the two calibrations
empirically.
\end{corollary}

%% W3-BEGIN (calibration bridge, exhibit TB; source:
%% agents/2026-08-20_sol_F1_report.md, calibration_bridge.csv and
%% frontier.csv; replaces the stale 2026-08-13 worked example):
Table~\ref{tab:bridge} evaluates both calibrations at the measured panel
inputs. The level calibration reproduces every adopted threshold on the theory
panels, and the frontier knee of Section~\ref{sec:scorecard} sits at the
calibrated value on three of the four frontiers.

\begin{table}
\begin{center}
\tbl{The level calibration reproduces the adopted thresholds, and the
frontier knee sits at the calibrated value on three of four panels. Inputs
measured per panel: span $N$, estimation step $k$ in observation periods,
excess kurtosis $\widehat{\kappa}$, and median within-cluster correlation
$\bar{\rho}$; thresholds from equations \eqref{eq:deltalevel} and
\eqref{eq:deltainnovation} at $z_{1-\alpha}=1$.}
{\begin{tabular}{lrrrrrrrl}
\toprule
panel & $N$ & $k$ & $\widehat{\kappa}$ & $\bar{\rho}$
 & $\delta^{*}_{\mathrm{lvl}}$ & $\delta^{*}_{\mathrm{inn}}$
 & adopted $\delta$ & knee \\
\colrule
equities, weekly & 156 & 4.33 & 2.12 & 0.62 & 0.087 & 0.028 & 0.0866
 & 0.0866 \\
futures, weekly & 156 & 4.33 & 1.61 & 0.68 & 0.069 & 0.023 & 0.0691
 & 0.0691 \\
funds, monthly & 36 & 1 & 0.84 & 0.80 & 0.083 & 0.027 & 0.0830 & 0.0830 \\
funds, quarterly & 12 & 1 & 1.29 & 0.79 & 0.161 & 0.089 & 0.1609 & 0.1609 \\
\colrule
U1 application, monthly & 36 & 1 & 2.12 & 0.64 & 0.175 & 0.058 & 0.0866
 & --- \\
U3 application, weekly & 156 & 4.33 & 1.61 & 0.68 & 0.069 & 0.023 & 0.0691
 & --- \\
\botrule
\end{tabular}}
\tabnote{The equity headline-window frontier places its knee at the fixed
bonus 0.05 rather than the calibrated 0.0866; Section~\ref{sec:scorecard}
reports the difference. Application rows state the operating cells of
Section~\ref{sec:applications}; their thresholds are transferred, not
recalibrated, as disclosed below.}
\label{tab:bridge}
\end{center}
\end{table}

The application rows make one transfer explicit. The adopted U1 threshold of
0.0866 is the weekly equity-panel calibration transferred to the monthly
span-36 application cell, whose own inputs imply a level threshold of 0.175.
We disclose the transfer rather than recalibrating, because the operating
points are frozen and the fidelity verdict of the adopted cell is quoted in
Section~\ref{sec:membership}.
%% W3-END

\subsection{Membership, lineage, and interpretability under smoothing}
\label{sec:membership}

%% TOPIC-CANDIDATE
The smoothed partitions must remain interpretable, and we measure
interpretability with the same instruments that measure stability. Per level
of an external taxonomy we report the adjusted Rand index of the partition
against the taxonomy. Per lineage track we report modal-taxonomy purity, the
share of track life under the modal label, and label churn per asset-year. The
fidelity band of equation~\eqref{eq:designproblem} is enforced as a symmetric
constraint: absolute taxonomy-ARI changes against the unsmoothed baseline of
at most $0.03$ at every level and a cluster-count change of at most $15\%$.

% Provenance: E3b corrected U1 fidelity verdicts,
% agents/2026-08-14_sol_E3b_report.md. Headline window rows.
On the equity panel, the calibrated bonus cuts raw membership churn from
$3.21$ to $0.36$ changes per asset-year (a reduction of 88.9\%) on the
headline window, while a quarterly-hold smoother reaches $1.33$ and a small
fixed bonus of $0.02$ reaches $1.27$. The reduction is not free: at the
calibrated bonus the sector-level taxonomy ARI moves by $-0.038$, which
breaches the $0.03$ band, while the small fixed bonuses stay inside the band.

%% W3-BEGIN (adopted-cell verdicts and interpretability; source:
%% agents/2026-08-20_sol_F3_report.md, adopted_cell_verdicts.csv and
%% interpretability.csv):
We score the two adopted application cells against their own unsmoothed
baselines, and both pass their bands. The monthly equity
cell at the transferred threshold changes taxonomy ARI by at most 0.015
against its own unsmoothed baseline, with a median cluster-count change of
0.12. The futures cell at 0.0691 changes taxonomy ARI by at most 0.009 with
no cluster-count change. Interpretability improves alongside stability.
Smoothing lowers annualised label churn from 1.16 to 0.64 changes per
asset-year on equities, from 0.54 to 0.10 on futures, and from 1.34 to 0.37
on funds. Modal-taxonomy purity rises on every panel, from 0.53 to 0.56, from
0.86 to 0.87, and from 0.69 to 0.73.
%% W3-END
\TODO{insert exhibit TC (churn and fidelity across panels and
configurations) and figure F1 (churn through time) from
finalisation/exhibits (F8). Band-verdict counts to re-trace: equities 6 of 8
configurations PASS on the full window and 5 of 8 on the headline window,
futures 8 of 8, funds 7 of 8 (F3 report).}

\subsection{Testable implications}
\label{sec:predictions}

The theory is falsifiable on the three estimation panels of
Section~\ref{sec:confirmation}. We state the implications as numbered
predictions. Each names its test statistic, and Section~\ref{sec:confirmation}
reports each with a bootstrap confidence interval or a permutation null.

\begin{prediction}[Churn concentrates at small margins]
\label{pred:margins}
Realised membership flips concentrate on assets whose estimated margins are
small relative to $\sigma_d(k)$, and the predicted churn
\eqref{eq:expectedchurn} matches realised churn across smoothing
configurations. Test: per-panel margin histograms, flip rates by margin
decile, and the cross-configuration correlation between predicted and realised
churn.
\end{prediction}

\begin{prediction}[The churn-fidelity frontier turns at the noise floor]
\label{pred:knee}
On the frontier traced by the bonus parameter $\delta$, churn falls steeply up
to the noise-floor calibration of Corollary~\ref{cor:delta} and fidelity
deteriorates beyond it. Test: sweep $\delta$ on a grid, overlay
$\delta^{*}_{\mathrm{lvl}}$ and $\delta^{*}_{\mathrm{inn}}$ computed from each
panel's span, frequency and kurtosis, and locate the frontier knee.
\end{prediction}

\begin{prediction}[Kurtosis raises the noise floor]
\label{pred:kurtosis}
Panels with fatter tails carry proportionally higher churn at a fixed
$\delta$, with the ordering given by $\sqrt{1+\kappa}$ estimated from the
data. Test: cross-panel comparison of realised churn against the
kurtosis-adjusted prediction, with $\kappa$ estimated from pooled returns and,
under the GARCH specification, from equation \eqref{eq:garchkappa}.
\end{prediction}

\begin{prediction}[Frequency scaling]
\label{pred:frequency}
% [TODO: OWNER (W2): restate P4 in its revised form per owner ruling 10 (the
% per-panel constant absorbed; the prediction concerns the k-dependence of
% churn ratios). The F2 re-evaluation is complete and its verdict is
% reported in Section 5; only the prediction text awaits the owner.]
Per-transition churn scales with the per-transition noise
$\sqrt{1-\lambda^{k}}$ when estimation dates move from monthly to quarterly on
the same panel, after absorbing a per-panel proportionality constant. The
optimal bonus $\delta^{*}$ moves with $k$ under the innovation calibration
\eqref{eq:deltainnovation} and stays fixed under the level calibration
\eqref{eq:deltalevel}. Test: re-score cached estimations at the coarser
schedule and compare both scalings against their predictions.
\end{prediction}

\begin{prediction}[Stabilisation is free at the risk-model level]
\label{pred:invariance}
At production regularisation, the fitted covariance is insensitive to the
partition, so smoothing changes the risk model by an order of magnitude less
than it changes the partition. Test: relative Frobenius distance and maximal
entry change of the fitted covariance, and the residual-diagonality
diagnostics, between smoothed and baseline runs on all three panels.
\end{prediction}

\begin{prediction}[Ergodic partition statistics]
\label{pred:ergodicity}
Partition summary statistics (cluster count, size entropy, consecutive-date
$\ARI$ and $\VI$) have stable subsample averages over the sample, consistent
with the stationarity and ergodicity of Lemma~\ref{lem:stationarity}. Test:
subsample means across sample halves and thirds against bootstrap bands, with
crisis windows reported separately.
\end{prediction}

\begin{prediction}[Turnover attribution in the momentum overlay]
\label{pred:turnover}
In a cluster-relative momentum portfolio, reassignment-driven turnover falls
monotonically in the smoothing strength while signal-driven turnover is
invariant, and net-of-cost performance is non-decreasing within the fidelity
band. Test: turnover decomposition against the prior-date partition
counterfactual, with block-bootstrap confidence intervals on the turnover and
net performance deltas.
\end{prediction}

% ----------------------------------------------------------------------------
\section{Theory confirmation on three estimation panels}
\label{sec:confirmation}
% ----------------------------------------------------------------------------

\subsection{Estimation panels and data}
\label{sec:panels}

%% TOPIC-CANDIDATE
We confirm the predictions on three cached estimation panels that differ in
asset class, cross-section size, frequency, and tail weight: a United States
equity panel (point-in-time index constituents, weekly observations, span
156), a global futures panel (95 contracts, weekly observations, span 156),
and a multi-asset fund panel (monthly observations at span 36 with a quarterly
sleeve at span 12, replicating a production configuration). The equity and
futures panels coincide at the asset level with two of the investment
universes of Section~\ref{sec:applications}. The fund panel appears only in
this section, and it carries the native monthly-versus-quarterly frequency
test of Prediction~\ref{pred:frequency}. Measured excess-kurtosis parameters
$\widehat{\kappa}$ are $2.12$ (equities, weekly), $1.61$ (futures, weekly),
$0.84$ (funds, monthly), and $1.29$ (funds, quarterly), so the fat-tails
ordering of Prediction~\ref{pred:kurtosis} is visible already in the inputs.
% Provenance: E1 data report kurtosis inputs,
% agents/ROADMAP_cluster_lineage_empirics.md E1 record.
\TODO{panel summary-statistics table (Part A; exhibit TA of the Sol TODO
file): one row per panel with asset count, frequency, span, sample range,
number of estimation dates, and the measured $\widehat{\kappa}$ values quoted
above,
which then move from this prose into the table. Every QF sample paper carries
such a table.}

\subsection{The theory scorecard}
\label{sec:scorecard}

%% W3-BEGIN (theory scorecard, exhibit TD; source:
%% agents/2026-08-21_sol_F5_report.md, theory_scorecard.csv; every value
%% carries its artifact path and hash in the scorecard file):
Table~\ref{tab:scorecard} reports one row per prediction with its test
statistic, measured value, and verdict. Intervals are 95\% percentile
intervals from a joint moving-block bootstrap with block length 6 and 2{,}000
draws.

\begin{table}
\begin{center}
\tbl{The scorecard supports four predictions outright, two in revised form,
and rejects one. Panels are equities / futures / funds; equity values quote
the full window with the headline window in parentheses where they differ.}
{\begin{tabular}{llll}
\toprule
prediction & statistic & measured value & verdict \\
\colrule
P1 margins & predicted--realised churn corr.
 & 0.863 (0.872) / 0.865 / 0.911 & supported \\
P1 simulation & flat-cut corr.; Ward corr., MAE
 & 0.980; 0.975, 0.033 & supported \\
P2 knee & frontier knee vs calibration
 & calibrated $\delta$ on 3 of 4 & supported \\
P3 kurtosis & absorbed constants $c$
 & 0.81 to 2.15 & supported, revised \\
P4 equality & mean realised$-$predicted ratio gap
 & 0.20 / 0.42 & rejected \\
P4 revised & cross-config corr.\ at fixed $c$
 & 0.992 (0.993) / 0.939 & supported, revised \\
P5 invariance & max resid.-diagonality; Frobenius
 & 0.014; 0.057 & supported \\
P6 ergodicity & consec.-ARI subsample ranges
 & above baseline on all panels & supported \\
P7 turnover & reassignment delta per rebalance
 & $-0.118$ / $-0.055$ / $-0.087$ & conj.\ rejected \\
\botrule
\end{tabular}}
\tabnote{P1 interval endpoints: equities full $[0.851, 0.871]$, equities
headline $[0.865, 0.878]$, futures $[0.814, 0.903]$, funds
$[0.885, 0.931]$. P7 intervals exclude zero on all three panels. Ward
carries no acceptance threshold: its measured accuracy is the result.}
\label{tab:scorecard}
\end{center}
\end{table}

The verdicts read as follows. Prediction~\ref{pred:margins} is supported: the
cross-configuration correlation between predicted and realised churn is
$0.86$ on the equity panel's full window and $0.87$ on its headline window,
$0.86$ on futures, and $0.91$ on funds, each with an interval excluding zero.
Prediction~\ref{pred:knee} is supported with one stated exception: the
maximum-curvature knee sits at the calibrated threshold on three of the four
frontiers, while the equity headline window places it at the fixed bonus
$0.05$. Prediction~\ref{pred:kurtosis} holds in the absorbed-constant form:
the per-panel baseline constant $c$ (realised churn over the
Gaussian-predicted sum) ranges from $0.81$ on the futures panel to $2.15$ on
the quarterly fund sleeve. The theory treats $c$ as a panel-level constant
rather than a prediction.
Prediction~\ref{pred:frequency} in its unscaled equality form is rejected:
realised annualised churn ratios exceed the predicted ratios on every equity
and futures configuration, with mean absolute gaps between $0.20$ and $0.42$.
Its revised form is supported as an ordering: with $c$ calibrated once on the
baseline configuration, the cross-configuration correlation between revised
predictions and realised quarterly churn is $0.99$ on equities and $0.94$ on
futures. Every quarterly level gap stays positive, and the fund panel is
descriptive under the different-sleeve ruling.
Prediction~\ref{pred:invariance} is supported: the maximal
smoothed-versus-baseline residual-diagonality change is $0.014$ across all
panels and configurations, an order of magnitude inside the $5\%$ guard.
Prediction~\ref{pred:ergodicity} is supported: subsample consecutive-ARI
ranges of the calibrated configurations sit above their baseline ranges on
every panel ($0.92$--$0.97$ against $0.55$--$0.67$ on equities,
$0.98$--$0.99$ against $0.89$--$0.95$ on futures, and $0.96$--$0.98$ against
$0.79$--$0.83$ on funds). The crisis windows preserve the ordering.
Prediction~\ref{pred:turnover} is rejected as a conjunction, and we report
which clause fails. Reassignment-driven turnover falls with smoothing on
every panel, with intervals excluding zero ($-0.118$, $-0.055$, and $-0.087$
changes per rebalance). The invariance clause fails on the fund panel, where
signal-driven turnover also moves ($-0.009$, interval excluding zero), and
the equity net-return point estimate is negative. We keep the mechanism and
drop the conjunction.
%% W3-END
\TODO{insert companion figures F2 (frontier knee with calibration overlays)
and F4 (margin histograms and flip rates by margin decile) from
finalisation/exhibits (F8); reference and interpret both here.}

\subsection{Simulation study}
\label{sec:simulation}

%% W3-BEGIN (simulation results; source: agents/2026-08-21_sol_F4_report.md,
%% simulation_results.csv and ward_verification.csv; 144 cells, 500 paths per
%% cell, seeds 20260817 + cell index recorded in the output):
We verify the flip approximation \eqref{eq:flipprob} on synthetic $G$-block
panels with GARCH(1,1) volatility across 144 design cells: dimensions 50 and
100, five and ten blocks, three separations, three volatility regimes, spans
36 and 156, two estimation steps, and four hysteresis arms per method. For the proved flat-cut case in the Gaussian
cells with separation at least $0.20$, the correlation between predicted and
realised flip probabilities is $0.98$, and realised churn is non-increasing
in the hysteresis strength in every one of the 144 cells. Ward linkage is
the verification case and carries no acceptance threshold. Its measured
correlation is $0.97$ with a mean absolute prediction error of $0.033$,
against $0.009$ for the flat cut on the same rows. The approximation
therefore transfers to Ward with a measured, roughly threefold, loss of
per-cell accuracy and no loss of ordering.
%% W3-END
\TODO{insert figure F3 (predicted versus realised flip probability, flat cut
and Ward) from finalisation/exhibits (F8); reference and interpret here.}

% ----------------------------------------------------------------------------
\section{The investment applications}
\label{sec:applications}
% ----------------------------------------------------------------------------

\subsection{Design and universes}
\label{sec:design}

%% TOPIC-CANDIDATE
The applications test one question in two settings: does replacing an external
peer definition with Rolling-Ward clusters improve a portfolio, holding
everything else fixed? The signal experiment changes only the peer set inside
which a momentum score is standardised. The risk experiment changes only the
grouping inside which a long-only risk allocation distributes its budgets.
Table~\ref{tab:universes} defines the three investment universes.

\begin{table}
\begin{center}
\tbl{Investment universes and frozen operating points.
\TODO{finalise caption with data sources, sample periods, and units per the
caption template of the Sol TODO file.}}
{\begin{tabular}{lccc}
\toprule
 & U1 equities & U2 funds & U3 futures \\
\colrule
universe & point-in-time MSCI US & fund catalogue, & continuous futures, \\
 & constituents & AUM $>$ USD 50m & 11 liquidity exclusions \\
clustering cell & monthly, span 36 & weekly, span 156 & weekly, span 156 \\
smoothing $\delta$ & 0.0866 & none & 0.0691 \\
signal & classic $12$m$-1$m & classic $12$m$-1$m & risk-adjusted momentum \\
portfolio & long-short & long-only & long-short \\
quantile & 25\% & 25\% & 25\% \\
sleeves & one cross-section & 55/35/10 E/FI/Rest & one cross-section \\
weighting & equal within side & equal within sleeve & inverse volatility \\
rebalancing & monthly & every two months & monthly \\
cost (one-way) & 10 bp & 20 bp & 10 bp \\
\botrule
\end{tabular}}
\label{tab:universes}
\end{center}
\end{table}
% Provenance: agents/2026-08-17_sol_signal_and_risk_model_pipeline_summary.md,
% selected signal specifications table.

Five controls apply throughout. Eligibility (membership, assets under
management, return history, and exclusions) uses only information available at
the decision date and is imposed before covariance estimation, clustering,
scoring, and construction. A weight decided at date $t$ is implemented with a
one-period lag. The backtests hold units between rebalances and compute
turnover and transaction costs from realised trades. Sharpe ratios use a zero
risk-free convention. The equal-weight basket of all eligible assets serves
only as the market reference for beta and alpha and never as a performance
yardstick. The common headline window is 2009-08-31 through 2026-06-30.

%% MECHANISM (candidate for owner rewrite): why peer-contained scoring.
In the signal experiment the same raw momentum panel
\citep{jegadeesh1993returns} is standardised in two ways: against the relevant
broad cross-section (the global control), and inside the asset's current
Rolling-Ward cluster (the treatment). For equities
a third leg standardises inside point-in-time sector classifications, which
tests the endogenous partition against the external taxonomy directly. A
cluster with ten or fewer members falls back to the global moments, which
prevents unstable small-group scores. Clusters receive no capital budgets: the
partition changes the score, not the allocation.

The sample vintages of the two experiments differ in two disclosed respects.
The risk-allocation runs of Section~\ref{sec:risk} were produced on an earlier
fund-eligibility rule (assets under management above USD 100m) and an earlier
futures exclusion set (seven contracts), while the final signal universes use
the USD 50m cutoff and eleven exclusions. The recorded samples are stated in
each exhibit and are not re-estimated. \TODO{owner sign-off on this
disclosure paragraph; ruling 5 fixed the analysis on the recorded samples
with bootstrap as the only addition.}

\subsection{Signal evidence: peer-contained momentum scoring}
\label{sec:signal}

%% TOPIC-CANDIDATE (re-based to the gated G0/F6 record in revision 3)
Cluster-contained scoring lowers the volatility of matched momentum
portfolios in all three universes, and the volatility reductions are the only
deltas whose bootstrap intervals exclude zero. Table~\ref{tab:signal} reports
the matched legs.

\begin{table}
\begin{center}
\tbl{Peer-contained momentum scoring against matched controls.
Annualised net returns after the stated one-way costs, annualised volatility,
zero-risk-free Sharpe ratios, and annualised one-way turnover, on
2009-08-31 through 2026-06-30.
\TODO{append the bootstrap confidence intervals (Part B) as companion
columns, not a separate table; finalise caption per the caption template of
the Sol TODO file.}}
{\begin{tabular}{llrrrr}
\toprule
universe & leg & net return & volatility & Sharpe & turnover \\
\colrule
% W3: U1 rows re-based to the gated G0 headline-window re-score
% (agents/2026-08-21_sol_G0_report.md, u1_windowed_performance.csv). The
% 2026-08-17 values are retired (lost source artifact; G0 reconciliation).
% U1 turnover awaits the exhibit-index trace (F8 table_existing_signal.csv).
U1 & cluster score & $-3.87\%$ & $9.17\%$ & $-0.38$ & \TODO{t/o} \\
U1 & sector score & $-3.52\%$ & $10.36\%$ & $-0.29$ & \TODO{t/o} \\
U1 & global score & $-3.94\%$ & $13.10\%$ & $-0.24$ & \TODO{t/o} \\
\colrule
U2 & cluster score & $5.48\%$ & $10.49\%$ & $0.56$ & $2.06\times$ \\
U2 & global score & $5.46\%$ & $11.17\%$ & $0.53$ & $1.76\times$ \\
\colrule
U3 & cluster score & $5.75\%$ & $8.18\%$ & $0.73$ & $2.27\times$ \\
U3 & global score & $6.35\%$ & $10.60\%$ & $0.63$ & $2.33\times$ \\
\botrule
\end{tabular}}
\label{tab:signal}
\end{center}
\end{table}
% Provenance: agents/2026-08-17_sol_signal_and_risk_model_pipeline_summary.md;
% runners: run_u1_classic_min_cluster_size_grid.py,
% run_u2_55_35_10_signal_grid.py, run_u3_rosaa_min10_short_span_sweep.py /
% run_u3_rosaa_min10_vol_normalized.py.

\TODO{figure F5 of the Sol TODO file: cumulative NAV of the cluster-score and
global-score legs per universe, one panel per universe, referenced and
interpreted here (QF precedent: the ESG paper's cumulative-factor-value
figure and BABB's risk-return profile).}

%% W3-BEGIN (comparison deltas and intervals; sources: U1 rows
%% agents/2026-08-21_sol_G0_report.md (u1_windowed_cis.csv), U2/U3 rows
%% agents/2026-08-20_sol_F6_report.md (signal_cis.csv); joint moving-block
%% bootstrap, block 6, 2,000 draws, 95\% percentile intervals):
Volatility falls in every universe against the matched global control: by 3.9
percentage points in U1 (interval $[-5.0, -2.8]$), 0.7 in U2
($[-1.2, -0.1]$), and 2.4 in U3 ($[-3.3, -1.6]$), and all three intervals
exclude zero. The return deltas are small and inconclusive: $+7$ basis points
in U1, $+2$ in U2, and $-60$ in U3, with every interval covering zero. The
Sharpe deltas are $-0.14$ in U1 ($[-0.34, +0.07]$), $+0.03$ in U2
($[-0.09, +0.14]$), and $+0.09$ in U3 ($[-0.13, +0.30]$), and none excludes
zero. Against the sector control in U1, the cluster score lowers volatility
by 1.2 percentage points ($[-2.1, -0.3]$) at a return delta of $-35$ basis
points and a Sharpe delta of $-0.09$, with only the volatility interval
excluding zero.
%% W3-END

%% MECHANISM (candidate for owner rewrite): what the result is and is not.
The supported claim is comparative and risk-focused. Standardising momentum
inside endogenous correlation clusters prevents heterogeneous instruments from
dominating one cross-sectional scale, and in the selected specifications this
reduces volatility in every universe while leaving the return and Sharpe
intervals covering zero. The evidence does not support an unconditional
return or Sharpe claim: the U1 long-short book loses money outright over this
sample under every scoring scheme, and the U3 cluster book trails its global
control on raw return. The U1 result in particular is a statement about
relative risk, and we report it as such.

\subsection{Risk evidence: cluster-budgeted long-only allocation}
\label{sec:risk}

%% TOPIC-CANDIDATE
The risk experiment uses no signal. At each estimation date the same
covariance matrix and asset set feed five allocation rules that differ only in
how they use the cluster structure: equal asset risk budgets (the flat ERC
control) \citep{maillard2010properties}, cluster risk budgets proportional to
$n_g^{\alpha}$ with $\alpha \in \{0.5, 0\}$ (square-root-size and
equal-cluster budgets, solved with the same production risk-budgeting solver),
hierarchical risk parity on the Rolling-Ward tree, and hierarchical risk
parity on canonical single linkage \citep{lopezdeprado2016building}.

\begin{table}
\begin{center}
\tbl{Signal-free long-only allocation across universes. Annualised net
returns after costs (10 bp one-way for U1 and U3, 20 bp for U2), volatility,
zero-risk-free Sharpe, and annualised one-way turnover.
\TODO{state the recorded sample vintages (U2 on the USD 100m rule; U3 on the
seven-exclusion universe) in the final caption; append bootstrap intervals
(Part B) as companion columns.}}
{\begin{tabular}{llrrrr}
\toprule
universe & method & net return & volatility & Sharpe & turnover \\
\colrule
% W3: the three G0-covered U1 legs re-based to the gated headline-window
% re-score (agents/2026-08-21_sol_G0_report.md). The two budget rows and all
% U1 turnover cells await the same-convention windowed values via the
% exhibit-index trace (F8 table_existing_risk.csv) -- re-scoring of frozen
% weights/NAVs only, no rerun. The retired full-range values must not return.
U1 & flat ERC & $8.79\%$ & $13.27\%$ & $0.71$ & \TODO{t/o} \\
U1 & cluster budgets, $\alpha=0.5$ & \TODO{tbd} & \TODO{tbd} & \TODO{tbd} & \TODO{t/o} \\
U1 & equal-cluster budgets & \TODO{tbd} & \TODO{tbd} & \TODO{tbd} & \TODO{t/o} \\
U1 & Rolling-Ward HRP & $9.00\%$ & $12.83\%$ & $0.74$ & \TODO{t/o} \\
U1 & single-linkage HRP & $8.86\%$ & $12.95\%$ & $0.72$ & \TODO{t/o} \\
\colrule
U2 & flat ERC & $0.94\%$ & $1.61\%$ & $0.59$ & $0.36\times$ \\
U2 & cluster budgets, $\alpha=0.5$ & $0.48\%$ & $1.30\%$ & $0.38$ & $0.39\times$ \\
U2 & equal-cluster budgets & $0.21\%$ & $1.12\%$ & $0.19$ & $0.44\times$ \\
U2 & Rolling-Ward HRP & $-0.17\%$ & $0.76\%$ & $-0.22$ & $0.52\times$ \\
U2 & single-linkage HRP & $-0.08\%$ & $0.53\%$ & $-0.15$ & $0.30\times$ \\
\colrule
U3 & flat ERC & $0.84\%$ & $1.67\%$ & $0.50$ & $0.24\times$ \\
U3 & cluster budgets, $\alpha=0.5$ & $0.99\%$ & $1.70\%$ & $0.58$ & $0.28\times$ \\
U3 & equal-cluster budgets & $1.23\%$ & $1.82\%$ & $0.67$ & $0.36\times$ \\
U3 & Rolling-Ward HRP & $0.03\%$ & $0.93\%$ & $0.04$ & $0.25\times$ \\
U3 & single-linkage HRP & $0.01\%$ & $0.94\%$ & $0.02$ & $0.41\times$ \\
\botrule
\end{tabular}}
\label{tab:risk}
\end{center}
\end{table}
% Provenance: agents/2026-08-16_sol_U{1,2,3}_hierarchical_risk_report.md and
% agents/2026-08-17_sol_signal_and_risk_model_pipeline_summary.md;
% runners: run_u1_hierarchical_risk.py, run_u2_hierarchical_risk.py,
% run_u3_hierarchical_risk.py, hierarchical_risk_allocations.py.

\begin{table}
\begin{center}
\tbl{Cluster-risk concentration under flat ERC and equal-cluster
budgets. Effective risk clusters is the inverse Herfindahl index of absolute
cluster-risk shares. \TODO{finalise caption per the caption template of the
Sol TODO file.}}
{\begin{tabular}{lrrrr}
\toprule
 & \multicolumn{2}{c}{effective risk clusters}
 & \multicolumn{2}{c}{largest cluster-risk share} \\
universe & flat ERC & equal-cluster & flat ERC & equal-cluster \\
\colrule
U1 & 39.7 & 60.4 & $6.3\%$ & $1.7\%$ \\
U2 & 5.9 & 14.0 & $32.5\%$ & $8.2\%$ \\
U3 & 9.4 & 16.2 & $21.7\%$ & $6.3\%$ \\
\botrule
\end{tabular}}
\label{tab:concentration}
\end{center}
\end{table}

\TODO{figure F6 of the Sol TODO file: the risk-concentration mechanism
through time (effective risk clusters or largest cluster-risk share, flat ERC
against equal-cluster budgets), referenced and interpreted here.}

%% W3-BEGIN (findings re-based to G0/F6 intervals; sources: U1 rows
%% agents/2026-08-21_sol_G0_report.md, U3 rows
%% agents/2026-08-20_sol_F6_report.md, risk_cis.csv):
Two findings are universe-dependent and one is universal. Rolling-Ward HRP is
the best performer in U1: against flat ERC it adds 21 basis points of net
return and 0.03 of Sharpe and it lowers volatility by 0.44 percentage points
($[-0.87, -0.08]$). The volatility interval is the only one excluding zero.
Rolling-Ward HRP also beats single-linkage HRP on return, volatility,
and Sharpe, with every interval covering zero. Equal-cluster budgeting is the
best performer in U3: it adds 39 basis points ($[+0.04, +0.77]$ percentage
points, excluding zero) and 0.17 of Sharpe ($[-0.05, +0.37]$) over flat ERC.
Volatility also rises by 0.15 percentage points ($[+0.01, +0.30]$, excluding
zero), so the U3 gain is return at slightly higher risk. In U2 flat ERC
remains superior, and the failure is informative: unconstrained HRP allocates
98.9\% to fixed-income funds with 2.5 effective assets, so its low volatility
comes from concentration in near-duplicate low-volatility instruments rather
than from diversification. The universal finding is the mechanism
(Table~\ref{tab:concentration}): equal-cluster budgeting raises the number of
effective risk clusters by 52\% to 136\% and cuts the largest cluster-risk
share in every universe. The cluster structure is an effective control layer
for distributing risk across endogenous peer groups. Whether that mechanical
diversification also raises realised Sharpe depends on universe composition.
%% W3-END

% ----------------------------------------------------------------------------
\section{Robustness and limitations}
\label{sec:robustness}
% ----------------------------------------------------------------------------

%% TOPIC-CANDIDATE
The robustness material comes from the selection grids that produced the
frozen operating points, and we disclose it as such: these grids are the
selection record, not independent confirmation.
%% W3-BEGIN (grid inventory; source: agents/2026-08-21_sol_F8_report.md,
%% robustness-grid consolidation, byte-identical copies):
Four grids are reported: the U2 eligibility grid across no cutoff, USD 50m,
and USD 100m, the U1 minimum-cluster-size grid around the adopted value of
10, the U3 short-span sweep around the adopted value of 3, and the covariance
frequency and span grids around the adopted monthly span-36 equity cell.
Appendix~\ref{app:robustness} reports each grid with its selection role
stated on every non-selected row.
%% W3-END
\TODO{insert appendix tables TE and TF from finalisation/exhibits (F8) and
add the one-paragraph reading of each grid from those tables.}

\label{sec:limitations}
%% TOPIC-CANDIDATE
Five limitations bound the claims. First, the fund catalogue is a
current-vintage snapshot, so the U2 results carry survivorship bias, and we
report them as comparative statements between legs run on the identical
sample. Second, the operating points were selected across grids on the full
sample, and the rolling-forward construction is the robustness argument: every
partition, score, and weight at date $t$ uses information up to $t$ only.
Third, the risk-allocation exhibits stand on earlier recorded eligibility
vintages, as disclosed in Section~\ref{sec:design}. Fourth, the futures
partitions were fitted before the final liquidity exclusions, and excluded
contracts carry zero weight rather than triggering a refit. Fifth, the
calibrated bonus breaches the taxonomy fidelity band on the weekly equity
cell, and the adopted monthly equity cell quotes its own band verdict in
Section~\ref{sec:membership}. \TODO{owner review of this section; the five
items mirror the honesty ledger of the status file and the wording is
owner-facing.}

% ----------------------------------------------------------------------------
\section{Conclusion}
\label{sec:conclusion}
% ----------------------------------------------------------------------------

%% TOPIC-CANDIDATE
We derived the distribution theory of Rolling-Ward partitions and turned the
cluster-smoothing parameter into a designed quantity. The stationary law is
inherited from the EWMA estimator, the increment law carries an effective
sample size of exactly $N$, and membership churn is a boundary-flip
phenomenon whose expected rate follows from the estimator's noise floor.
%% W3-BEGIN (quantified contributions; sources: F5 scorecard, F1 bridge,
%% G0/F6 intervals, F4 simulation):
The evidence supports four of the seven predictions outright and two in the
revised absorbed-constant form, and it rejects the turnover conjunction while
confirming its reassignment mechanism with intervals excluding zero on all
three panels. The level calibration reproduces the adopted thresholds and
locates the frontier knee on three of four panels, and the simulated flip
approximation reaches a predicted-realised correlation of $0.98$ for the
proved flat cut and $0.97$ for Ward. The calibrated hysteresis threshold cuts
membership churn by up to 89\% at a bounded and measured fidelity cost. In
the applications, the stabilised clusters lower matched signal-portfolio
volatility by 0.7 to 3.9 percentage points with every interval excluding
zero, and cluster risk budgets raise the number of effective risk clusters by
52\% to 136\% in the allocation application.
%% W3-END

% ----------------------------------------------------------------------------
% Notation table: [TODO: add after the model sections stabilise, per the
% notation-budget rule. Acronym audit measured 2026-08-17 (body, comments
% excluded): HRP 12, ERC 12, EWMA 9, GARCH 8, ARI 4, VI 3 -- ARI and VI
% survive as math operators carried by the notation table; AUM now spelled
% out in prose (table cells keep the short form, defined in captions); MCF
% no longer appears in the body and is struck from this audit; FF6/MATF do
% not appear and are struck.]
% ----------------------------------------------------------------------------

% ----------------------------------------------------------------------------
% Taylor & Francis back matter, in the order and with the statements the QF
% Instructions for Authors (17 July 2026) require: acknowledgements,
% disclosure statement, declaration of generative AI use (REQUIRED),
% funding (declare "no funding was obtained" if none), data availability.
% ----------------------------------------------------------------------------

\section*{Acknowledgements}
\TODO{acknowledgements.}

\section*{Disclosure statement}
The authors report there are no competing interests to declare.
\TODO{owner to confirm, or list interests; wording above is the journal's
suggested form.}

\section*{Declaration of generative AI use}
Generative AI (Anthropic Claude) was used in the preparation of this
manuscript under the direction of the authors: drafting and editing of text
from author-written outlines, central claims, and rulings, and consistency
checks of quoted numbers against the replication outputs. All results,
citations, and final wording were reviewed and approved by the authors, who
take full responsibility for the content of the article. \TODO{owner to
confirm this wording against the T\&F AI-use guidance referenced in the
Instructions for Authors; this statement is mandatory (state explicitly if
generative AI had not been used).}

\section*{Funding}
\TODO{funding statement; if none, use the journal's suggested form: no
funding was obtained for this work.}

\section*{Data availability statement}
%% W3-BEGIN (F10 draft statement from
%% agents/2026-08-21_sol_F10_replication_statement.md, with its two
%% semicolon sentences split per the house rule; owner sign-off under W2
%% before submission):
The empirical analysis uses licensed MSCI constituent data, Bloomberg-derived
fund and futures data, and factor returns from the Ken French Data Library.
The Ken French series are the only freely obtainable data input. The licensed
MSCI and Bloomberg inputs cannot be redistributed, and researchers seeking to
reproduce the empirical results must obtain the corresponding data licences
and construct the documented point-in-time panels. The \texttt{factorlasso}
implementation of Rolling-Ward clustering, smoothing calibration, and
cluster-lineage analytics is publicly available. The paper-specific
replication harness is available from the authors on request and provides the
versioned scripts, frozen specifications, provenance records, deterministic
seeds, and exhibit-building instructions needed to reproduce the analysis
once the licensed inputs have been supplied. The request package does not
contain or purport to provide the proprietary source data.
%% W3-END
\TODO{owner sign-off (W2). The journal supports Open Science Badges (COS):
the Open Materials badge fits the code-public/data-proprietary split, and a
deposited replication archive with a DOI supports both this statement and
the badge application, which is made on acceptance.}

\bibliographystyle{rQUF}
\bibliography{cluster_lineage_refs}

% ----------------------------------------------------------------------------
% Appendices follow the reference list in rQUF style (rQUFguide sec. 4.11).
\appendices
% ----------------------------------------------------------------------------
\section{Proofs}
\label{app:proofs}

% Proofs moved from the body in revision 2 (E4). Content unchanged.

\begin{proof}[Proof of Lemma~\ref{lem:stationarity}]
The series \eqref{eq:stationarysol} converges almost surely under the
logarithmic moment condition, because $\lambda^j$ decays geometrically while
$\log^{+}\|r_{t-j} r_{t-j}^{\top}\|$ grows sub-linearly along stationary
sequences. The recursion is an iterated random function with a constant
contraction factor $\lambda < 1$, so the backward iteration argument of
\citet{diaconis1999iterated} gives uniqueness of the stationary law and
ergodicity. Measurable functions of stationary ergodic processes are stationary
and ergodic, which covers $Q_t$ and any partition statistic.
\end{proof}

\begin{proof}[Proof of Proposition~\ref{prop:increments}]
Part (i): $\vech(\widehat{\Sigma}_t)$ is a linear filter of the i.i.d.\
sequence $\vech(r_t r_t^{\top})$ with coefficients $(1-\lambda)\lambda^{j}$, so
$\gamma(0) = (1-\lambda)^2 V / (1-\lambda^2) = V(1-\lambda)/(1+\lambda)$, and
substituting $\lambda = 1 - 2/(N+1)$ gives $V/N$. The autocovariance follows
from the same filter representation. Part (ii) is the standard identity for
stationary AR(1)-type autocovariances. Part (iii) applies the delta method for
correlations to the fourth-moment structure of elliptical laws, for which the
asymptotic variance of $\widehat{\rho}$ is $(1+\kappa)(1-\rho^2)^2$ per
effective observation. Part (iv) rearranges the unconditional kurtosis of the
Gaussian GARCH(1,1), $K = 3\,(1-(\alpha+\beta)^2)/(1-(\alpha+\beta)^2-2\alpha^2)$,
through $K = 3(1+\kappa_{\Theta})$.
\end{proof}

\begin{proof}[Proof sketch of Proposition~\ref{prop:flip}]
A flip requires the noisy margin to cross zero against the hysteresis band,
which is a Gaussian tail event for the difference of two estimated mean
distances, whence the factor $\sqrt{2}$. Summing marginal flip probabilities
gives the expected count without any independence requirement, because
expectation is linear. The Gaussian approximation for $\widehat{m}_{i,t}$
follows from Proposition~\ref{prop:increments} and the delta method under
Assumption~\ref{ass:block}, where merge heights concentrate and the population
tree is deterministic.
\end{proof}

\section{Selection grids}
\label{app:robustness}

\TODO{Part B / Phase C: the four selection-grid tables of
Section~\ref{sec:robustness}, exhibits TE and TF of the Sol TODO file (grids
consolidated two per table). Appendix placement per the QF-baseline exhibit
budget; each table labelled with its selection role.}

\end{document}
